\newpage
\section{Abordagem estatística da lei do valor}
\subsection{MBA}

\begin{quote}
\textit{For prices at which commodities are exchanged to approximately correspond to their values, nothing more is necessary than 1) for the exchange of the various commodities to cease being purely accidental or only occasional; 2) so far as direct exchange of commodities is concerned, for these commodities to be produced on both sides in approximately sufficient quantities to meet mutual requirements, something learned from mutual experience in trading and therefore a natural outgrowth of continued trading; and 3) so far as selling is concerned, for no natural or artificial monopoly to enable either of the contracting sides to sell commodities above their value or to compel them to undersell. By accidental monopoly we mean a monopoly which a buyer or seller acquires through an accidental state of supply or demand.}

\textit{The assumption that the commodities of the various spheres of production are sold at their value merely implies, of course, that their value is the centre of gravity around which their prices fluctuate, and their continual rises and drops tend to equalise.}

(Marx. Capital, Vol 3. 1971: 178)

\end{quote}

\textbf{Summary generated by alphavix}

This paper investigates Marx's theory of the law of value using a dynamic computational model. Here's a summary of its key findings and contributions:

The paper aims to verify Marx's claim that prices gravitate towards labor values in a simple commodity economy, where capitalist investment and profit are absent. This is a crucial aspect of Marx's theory, as the law of value forms a fundamental building block. The author employs an agent-based computational model to simulate a simple commodity economy with  economic actors and  commodity types. The model incorporates production, consumption, exchange, and a market mechanism with subjective pricing, along with a rule for actors to switch production sectors.

The main conclusions drawn from the computational model are:

\begin{itemize}
    \item Labor value as an attractor for market price: The model demonstrates that the labor value of a commodity acts as an attractor for its market price, meaning that market prices tend to converge towards the labor values over time.
    \item Market prices as error signals: Market prices serve as error signals that dynamically allocate social labor between different production sectors.
    \item Deviations between prices and values indicate misallocations of labor. Efficient allocation of social labor: The tendency of prices to approach labor values is the monetary manifestation of a simple commodity economy efficiently allocating social labor. This suggests that the market mechanism, even with subjective pricing, can lead to an efficient division of labor.
    \item Emergence of Marx's law of value: The model shows that Marx's law of value can spontaneously emerge from numerous local exchanges, operating through money flows that constrain actors' subjective price evaluations. This happens "behind the backs" of individual actors, without their explicit awareness or intent.
\end{itemize}
The paper also introduces the concept of the Monetary Expression of Labour-time (MELT), which quantifies the constant of proportionality between mean prices and labor values. It highlights MELT as a measure of how much labor-time money represents and its role in connecting non-market phenomena (production times) with market phenomena (prices). Overall, the research supports the theoretical integrity of Marx's law of value in the context of a simple commodity economy, providing a dynamic analysis that complements traditional static models. It suggests that the market, through price signals and money flows, effectively coordinates the allocation of social labor, leading to an equilibrium where prices reflect the underlying labor values.


\textbf{Statistical Approximation of the Law of Value}

This model is presented in chapter 9 of the book \textit{Classical Econophysics} and also as the paper: \href{https://doi.org/10.1080/09538250701661889}{\textit{The Emergence of the Law of Value in a Dynamic Simple Commodity Economy}}.

 First, we must clearly establish that there is a distinction between value and price:

\begin{itemize}
    \item \textbf{Value}: determined by the prevailing technical conditions of production and measured by the socially necessary labor time to produce it.
    \item \textbf{Price}: the amount of money that the commodity yields in the market.
\end{itemize}
In a theoretical simplification of capitalism (a "simple commodity economy"), prices tend to values, and this is what we want to demonstrate.

\subsubsection{The Model}

The essential characteristics of the model are:

\begin{itemize}
    \item It is composed of $N$ workers (identified by an integer $i$ between 1 and $N$).  
    \item There are $L$ commodities (identified by an integer $j$ between 1 and $L$).  
    \item It has a total and constant amount of money $M=\sum_{i}^{N}m_{i}$.  
\item Each worker produces one commodity at a time.  
\item Each commodity is simple: it does not require another commodity to be produced and can be produced by a single worker, and all workers produce the same commodity with the same efficiency.  
\item The agents produce to satisfy their needs.
\end{itemize}


\textbf{Production Rule $P_{1}$}:

\begin{itemize}
    \item  At the beginning of the simulation, each agent $i$ has a vector $\boldsymbol{e}_{i}=\boldsymbol{0}$ that indicates the quantity of commodities agent $i$ possesses.  
\item The commodity $j$ being produced by agent $i$ is given by $A(i)=j$.  We can think of a vector $\boldsymbol{A}$ that tells us what each agent $i$ is producing at the moment.  
\item Each commodity requires $l_{j}$ steps to be produced.    That is, at each step the agent produces $L=1/l_{j}$ units of the commodity.   We define a vector $\boldsymbol{l}=\left(1/l_{1},\dots,1/l_{L}\right)=\left(L_{1},\dots,L_{L}\right)$.  
\end{itemize}

Thus, agent $i$ generates one unit of commodity $A(i)$ every $l_{A(i)}$ steps, and as a consequence the element of the vector $e_{i}[A(i)]$ is incremented by one unit.


\textbf{Consumption Rule $C_{1}$}

\begin{itemize}
    \item  All agents have the same consumption desire given by a global vector  
  $\boldsymbol{c}=\left(1/c_{1},\dots,1/c_{L}\right)=\left(C_{1},\dots,C_{L}\right)$.  
\item Each agent $i$ has a consumption deficit vector initialized as $\boldsymbol{d}_{i}=0$.  
\item Analogous to production, every $c_{j}$ steps the element $\boldsymbol{d}_{i}[j]$ is incremented by $1$,    that is, at each step the agent’s desire increases by $C_{j}=1/c_{j}$ unit.  
\end{itemize}

Thus, at each step, agent $i$ consumes a quantity of commodities given by the vector  $\boldsymbol{o}_{i}=\text{min}\left(\boldsymbol{e}_{i},\boldsymbol{d}_{i}\right)$.  For example, for commodity $j=1$:  
\begin{itemize}
    \item  If the agent has none of it, $e_{i}[1]=0$, then they cannot consume.  
\item If they have no desire, $d_{i}[1]=0$, they also will not consume.  
\item If both values are nonzero and the agent has more than they desire, $e_{i}[1]>d_{i}[1]$,    then they consume only what they desire, $d_{i}[1]$.  
\item If the agent has fewer goods than they desire, $e_{i}[1]<d_{i}[1]$,    they consume all they have, $e_{i}[1]$.  
\end{itemize}

In all cases, consumption is given by the smaller value:  
$o_{i}[1]=\text{min}\left(e_{i}[1],d_{i}[1]\right)$.  Evidently, this consumed amount must be deducted from both the deficit and the commodities in possession.  Thus, the vectors are updated as:  $\boldsymbol{e}'_{i}=\boldsymbol{e}_{i}-\boldsymbol{o}_{i}$ and $\boldsymbol{d}'_{i}=\boldsymbol{d}_{i}-\boldsymbol{o}_{i}$.


\textbf{Reproduction Coefficient $\eta=\sum\frac{l_{j}}{c_{j}}=\sum\frac{C_{j}}{L_{j}}$:  }
\begin{itemize}
    \item  $\eta=1$ means that production equals consumption: the simulation will adopt this condition.  
\item $\eta>1$ means the economy is permanently in deficit.  
\item $\eta<1$ means the economy permanently produces a surplus.  

\end{itemize}

Since under no circumstance do we have $c_{j}<0$ or $l_{j}<0$, to avoid $\eta>1$ we must ensure that no term is $\frac{l_{j}}{c_{j}}>1$.  If there is more than one commodity, then we must be even more restrictive and require $\frac{l_{j}}{c_{j}}<1$.
In other words, it must take fewer steps to produce a commodity than to desire it.  This is a necessary condition for stability, since each agent produces only one commodity at a time but consumes all.  For example, if we have two commodities, two agents, and each produces one commodity with values $l_{j}=1$ and $c_{j}=2$,  then $\eta=0.5+0.5=1$.  


\textbf{Price Rule $O_{1}$}

\begin{itemize}
    \item The price of commodity $j$ according to agent $i$ is a value $p_{j}^{i}$  
  that is randomly drawn from the interval $\left[0,m_{i}\right]$.
\end{itemize}


 \textbf{Market Rule $M_{1}$}

The market for a given commodity is considered "cleared" when there are no more buyers or sellers for that commodity.  In other words, if it is not cleared, it means there are still buyers and sellers.  We begin with a set $C$ of commodities that have not yet been cleared in the market:
\begin{enumerate}
    \item  Randomly select a commodity $j$ from the set $C$.

\item Randomly select a seller agent $s$ from the set of potential sellers:     that is, agents who have more of commodity $j$ than they wish to consume, $e_{i}[j]>d_{i}[j]$.

\item Randomly select a buyer agent $b$ from the set of potential buyers:     that is, agents who have fewer units of commodity $j$ than they wish to consume, $e_{i}[j]<d_{i}[j]$.

\item If there are no potential buyers or sellers, remove commodity $j$ from $C$.     If there are, apply the exchange rule $\boldsymbol{E}_{1}$.

\item Repeat the previous steps until all commodities are cleared.
\end{enumerate}


\textbf{Exchange Rule $E_{1}$}

The previous rule identifies buyers and sellers to carry out the conditional exchange defined here.

\begin{itemize}
    \item  Once we have a buyer $b$ and a seller $s$ who estimate the prices $p_{j}^{b}$ and $p_{j}^{s}$ for commodity $j$ by $O_1$ rule,    a transaction price is drawn from the discrete interval $\left[p_{j}^{b},p_{j}^{s}\right]$.  

\item If the buyer has enough money, the exchange takes place:  
  the buyer loses money and gains one unit of the commodity, while the seller gains money and loses one unit of the commodity.
\end{itemize}


\textbf{Sector Rule $S_{1}$}

\begin{itemize}
    \item After a fixed amount of time, considering we are at period $n$, each agent calculates an error vector  $\left\Vert \boldsymbol{d}_{i}^{n}\right\Vert$ and compares it with the same vector calculated in the previous period $\left\Vert \boldsymbol{d}_{i}^{n-1}\right\Vert$.  If the error has increased, $\left\Vert \boldsymbol{d}_{i}^{n}\right\Vert > \left\Vert \boldsymbol{d}_{i}^{n-1}\right\Vert$,  then the agent randomly switches the commodity being produced.


\end{itemize}


 \textbf{Simulation Rule $R_{1}$}

The entire cycle of production, consumption, exchange, and reallocation in production follows this rule.  Initially, we construct $\boldsymbol{l}$ and $\boldsymbol{c}$ such that $\eta=1$, and allocate $M/N$ among all agents. Then:

\begin{itemize}
    \item Increase the simulation time by one step.  
\item Invoke rule $P_{1}$ for each agent.  
\item Invoke rule $C_{1}$ for each agent.  
\item Invoke the market rule $M_{1}$.  
\item Invoke rule $S_{1}$ for each agent.  
\item Repeat.  
\end{itemize}


That is:

$$ SCE=\left\{ R_{1},P_{1},C_{1},O_{1}\left\{ M_{1},E_{1}\right\},S_{1}\right\} $$


The parameters are:

\begin{itemize}
    \item $N$ — number of agents.  
\item $L$ — number of commodities.  
\item $M$ — amount of money in the simulation.  
\item $R$ — the maximum possible consumption time, used to constrain the construction of the vectors $\boldsymbol{l}$ and $\boldsymbol{c}$.  

\item $C$ — a constant multiple of $R$, and a parameter that defines the duration of the period between applications of the sector-switching rule $S_{1}$.

\end{itemize}

\subsubsection{Division of Labor}

\textbf{Definition 1}: A division of labor is efficient when, for each commodity,  the number of commodities being produced equals the demand.   Thus, the total quantity of commodity $j$ demanded per step is $NC_{j}$.  If we have a fraction $\alpha N$ of agents producing commodity $j$,   then $a_{j}NL_{j}$ units are produced per unit of time.  

To achieve an efficient division of labor:


\begin{equation}
a_{j}NL_{j}=NC_{j}\quad\rightarrow \quad a_{j} =\frac{C_{j}}{L_{j}}=\frac{l_{j}}{c_{j}}
\end{equation}

In other words, $a_{j}$ represents an efficient division of labor.

 
\subsubsection{ Objective Prices}

We begin with two definitions:  

\begin{itemize}
    \item The average price of commodity $j$ is given by $\left\langle p_{j}\right\rangle$,    and we have the vector $\boldsymbol{p}=\left(\left\langle p_{1}\right\rangle ,\dots,\left\langle p_{L}\right\rangle \right)$.  
\item The values, i.e., the amount of time embedded in the production of the commodities, are $\boldsymbol{v}=\left(l_{1},\dots,l_{L}\right)$.  

\end{itemize}

Thus, if the price gravitates around value, if it is proportional to value, we can write:

\begin{equation}
\boldsymbol{p} \approx \lambda \boldsymbol{v}
\end{equation}

where $\lambda$ has units of 'money per unit of labor time'. Since our economy is simple, we can then define the Monetary Expression of Labor Time (MELT)  as the ratio between the measure of the total quantity of commodities exchanged over a time interval at current prices  and the productive labor expended in producing these commodities.  

In other words, MELT is simply:



\begin{equation}
\lambda = \frac{\gamma M}{\sum l_{i}v_{i}}
\end{equation}



Analyzing:  
\begin{itemize}
    \item The denominator is the labor expended on the commodities that are exchanged over a given time interval.    Here, $v_{i}$ is the average exchange rate of commodity $j$.  Thus, we have a summation where each term corresponds to the labor involved in the exchanges of each type of commodity.  

\item The numerator is the amount of money exchanged over a given time interval through commodity exchange.    Here, $\gamma$ is simply the proportion of money exchanged per unit of time, so $\gamma M$ is the velocity of money,    i.e., how much money is exchanged within the time interval under consideration.  

\end{itemize}

In summary, we have:

\begin{equation}
\lambda = \frac{\text{money exchanged per unit of time}}{\text{labor time exchanged per unit of time}} = \frac{\gamma M}{\sum l_{i}v_{i}}
\end{equation}

Evidently, we only focus on commodities that are exchanged, not on production as a whole; we are analyzing the market. It is worth noting that in this model commodities have a price only at the moment of exchange,  while their value is given by the technical production characteristics of the entire society.

\subsubsection{Simulation} 

Neither the book nor the paper provide the exact implementation code of the model, so I am trying to be as faithful as possible.  However, after several weeks working on this model, some comments need to be made:

\textbf{Comment 1}: At no point is it specified whether money is integer or not.  It is only mentioned that the interval from which prices are drawn is discrete,  which is a requirement for any computational model, no matter how small the intervals are.  

The exchange rule can make execution excessively slow if continuous (we will use "continuous" to indicate that we are using variables such as float or double)  money is used,  because the model does not advance until buyers and sellers reach an agreement to clear the market.  However, if a buyer has wealth close to $m_a \approx 0$, the commodity price defined by rule $O_1$  needs to be close to $0$ for the agent to be able to buy.  The lower the probability of this happening, the higher the wealth of the seller.  

In the appendix, the book mentions implementing a limit on the number of times each agent can enter the market,  but it does not detail exactly how this is done.  I opted to use decimal money; tests with integers produced faster but less precise results.  I also limited the agents' visits to the market, regardless of success or failure, or which commodity was involved. In my tests, increasing the number of market visits up to the tested limit made the model more precise but slower.  However, I am not sure whether allowing unlimited attempts necessarily improves the results.  It would require testing to verify whether failing to obtain the desired commodity due to insufficient funds is an important part of the system's dynamics.


\textbf{Comment 2}: In the implementation of rule $S_1$ for sector switching, it is not clear how the vector $\boldsymbol{d}_{i}^{(n)}$ is constructed. Is it simply the vector $\boldsymbol{d}_{i}$ at the last step of the period, or is it some distinct vector? I opted for this simpler implementation mentioned.

Also, in this same rule, it is stated that if there was an error increment, the agent should move to a \textit{new sector}* but it says that the new sector is drawn uniformly among the $L$ sectors. It is ambiguous whether the agent can end up remaining in the same sector. I implemented it so that the agent is forced to switch sectors. Some tests also suggest that this dynamic of agents switching between sectors plays an important role in distributing money among agents and sectors within the system, even when the system is in statistical equilibrium.

\textbf{Comment 3}: It was not clear to me what happens if, when applying rule $E_1$, the seller's offer $p_j^{(s)}$ is greater than the buyer's offer $p_j^{(b)}$. I assume that we always take both the buyer's and seller's offers and draw the final price between the smaller and the larger of the two, regardless of who is making the offer.  

In the application of rule $O_1$, it is allowed for a commodity to have a price $p=0$. Depending on the total amount of money, it is highly possible that a commodity could be exchanged at no price if we work with integers, which seems undesirable to me.


\textbf{Comment 4}: To conclude, I must admit that I had difficulty obtaining the results presented for $L>2$. I initially tried to implement the code in Python, but then migrated to C\# in search of higher execution speed. My results show a clear tendency for prices to be proportional to value, but not always with as strong a correlation as presented in both the book and the paper.

I believe that, besides some possible implementation differences, there are certain distributions of production rates and deficits for which the model can exhibit a higher correlation rate, and some parameter choices can make the model more accurate while at the same time increasing computational cost. In my tests, as mentioned, increasing the number of market visits per agent had a positive effect on the final result. Increasing the number of agents itself also helped, as well as using continuous money instead of integer. All these decisions improve the result at the expense of higher computational cost.

Increasing the data capture window—for example, averaging over 50,000 steps instead of 5,000—also had a positive effect, again at the cost of increasing the time required for the simulation to run. This is why I ultimately implemented the code in C\# (I apologize for the lack of comments. I spent so much time trying to reproduce this model that I got exhausted. In the future, I may revisit it and organize the code properly).


However, I must warn that I still believe there may be some subtle differences compared to the original proposal due to the difficulty of faithfully reproducing the results. The correlation originally reported, even for $L=5$, was around 9.8 over 10 runs, and in one case even reached 1. Nonetheless, I believe that this version of the model is already sufficient to investigate the issue. Suggestions for improvement are welcome.  

One of the main problems I notice is that the commodity with the longest production time seems to struggle to achieve the highest price. If we have $L=5$, the other four commodities follow a very strong linear trend, while this last commodity has difficulty showing the expected results. Additionally, this commodity also exhibits an extremely low number of transactions, which I believe is a result of the deficit accumulated in previous steps and the fact that the model prioritizes consumption over exchange. Thus, I assume (though further investigation is needed) that agents who start working in this sector largely end up producing more for consumption than for exchange, even more so than in the other sectors.

Below is a result for $L=5$, with $N=200$ agents, $M=100 \cdot N$, and $L \cdot 500$ market visits per step for each agent, data obtained over 50,000 steps (50,000-100,000):


\begin{verbatim}
Commodity 0 | Price: 38.33  | Exchanges: 1975   | Workers: 20 
Commodity 1 | Price: 116.84 | Exchanges: 154118 | Workers: 30 
Commodity 2 | Price: 553.42 | Exchanges: 22738  | Workers: 40 
Commodity 3 | Price: 747.60 | Exchanges: 2207   | Workers: 50 
Commodity 4 | Price: 541.66 | Exchanges: 148    | Workers: 60 
\end{verbatim}



The price and worker values are averages. We can note that the distribution of workers in each sector follows the efficient distribution, and that the price increases with value, except for the highest-value commodity. However, it has an exchange volume 13 times lower than the second-lowest traded commodity. One commodity was exchanged roughly every 500 steps.  If we take into account the number of transactions per agent, the figure becomes even lower. The number of exchanges per agent, for each commodity, is: $[98.75, 5137.26, 568.45, 44.14, 2.46]$.

It is worth remembering that there is a passage in the book discussing the requirements for the law of value to operate:

\begin{quote}
[if] the probability that a seller of j will find a buyer in the marketplace is low; hence exchange becomes occasional, failing a requirement for the law of value to operate.

\end{quote}

I believe this is what happens in this model. Despite everything, the relationship between value and price remains clear: If you increase the average production time of a commodity, the average price at which this commodity will be traded tends to increase. The correlation considering all commodities is 0.84, and disregarding the highest-value commodity it is 0.97. The production time for each commodity is $[1,1.5,2,2.5,3]$, and the time required to increase a unit of deficit for all commodities is $10$.

 Due to the statistical nature of the model, when there are more exchanges, the law of value is respected. For almost the same case, but now allowing only $L \cdot 100$ market visits per agent, and taking the averages between steps 950,000 and 1,000,000, we have:



\begin{verbatim}
Commodity 0 | Price: 36.69  | Exchanges: 7151   | Workers: 20
Commodity 1 | Price: 101.88 | Exchanges: 139728 | Workers: 30
Commodity 2 | Price: 547.72 | Exchanges: 17750  | Workers: 40
Commodity 3 | Price: 666.68 | Exchanges: 1226   | Workers: 50
Commodity 4 | Price: 891.73 | Exchanges: 142    | Workers: 60
\end{verbatim}

A correlation of 0.97.  Another indication that our model is not entirely faithful to the book is that more than half of the commodities have an average price up to 5 times the per capita wealth. In the graph presented in the book (or paper), no wealth exceeds the per capita wealth. However, I believe we are still able to demonstrate the relationship between price and value.

It should also be considered that, although this result is proportional, it is not necessarily directly proportional. That is, we have a kind of base price $b$ for the commodity, to which more price is added according to the labor time required to produce it: $\left\langle p_{j}\right\rangle = \lambda l_{j} + b$. Alternatively, we can think of it as a kind of added value $n_{j} = \frac{b}{\lambda}$ in the commodity, which is not directly defined but emerges from the dynamics of the model. Like an extra time it takes on average to be exchanged after being produced: $\left\langle p_{j}\right\rangle = \lambda \left(l_{j} + n_{j}\right)$.  

We can try to describe the relationship between production and price with a line given by:

We could spend time interpreting what this increase (or decrease) in the production times of each commodity means, or the existence of a base price for the commodity, but the most relevant point, for me, is the undeniable linear relationship between price and value.
Furthermore, other modifications of this model can be explored with the aim of correcting and/or improving the results. Examples include: testing different rules for constructing the deficit vector to apply to the $S_1$ rule, or even an entirely different rule for $S_1$; alternative ways to limit the number of market visits per step; limiting each agent’s consumption per step; and so on. To conclude, we can analyze a simple special case with two commodities, in which it is easier to achieve equilibrium and we can efficiently allocate workers between the two sectors without major dynamic problems by disabling the switch rule.

This line is not the result of interpolating the points, but rather a line $y=x$ constructed to pass through the origin.Evidently, there is no merit in constructing a line that passes through two points, but nothing requires the line to pass through the origin. This is a result of the model.


\textbf{Extra}:

A possible algorithm for constructing the vectors $l$ and $c$ is as follows:   First, generate two random vectors $\boldsymbol{l} = \left(L_{1}, \dots, L_{L}\right)$ and $\boldsymbol{u} = \left(u_{1}, \dots, u_{L}\right)$. Compute the sum $\sum_{i}^{L} \frac{u_{j}}{L_{j}} = F$, then divide both sides by $F$:


\begin{equation}
\sum_{i}^{L}\frac{1}{F}\frac{u_{j}}{L_{j}}=\frac{F}{F}\rightarrow\sum_{i}^{L}\frac{u_{j}/F}{L_{j}}=1=\sum_{i}^{L}\frac{C_{j}}{L_{j}}
\end{equation}

That is, the elements of the vector $\boldsymbol{C}$ are $C_{j} = \frac{u_{j}}{F}$, or equivalently, $\boldsymbol{C} = \frac{\boldsymbol{u}}{F}$.
\subsection{MBE}


This model is presented in chapter 9 of the book \textit{Classical Econophysics} and also as the paper: \href{https://doi.org/10.1080/09538250701661889}{\textit{The Emergence of the Law of Value in a Dynamic Simple Commodity Economy}}.

\textbf{Summary generated by alphavix}

The paper aims to verify Marx's claim that prices gravitate towards labor values in a simple commodity economy, where capitalist investment and profit are absent. The mathematical modeling focuses on a system of ordinary differential equations to describe the dynamics of the simple commodity economy. This approach is intended to provide an intuitive explanation of the causal features observed in the computational model, rather than definitive proofs or a precise stochastic theory. Here's a breakdown of the key mathematical components:
\begin{itemize}
    \item Labour Equation: This equation describes how the allocation of labor to different production sectors changes over time. It's driven by a "sectoral income error" ($\Phi_j(t)$), which represents the difference between a sector's actual income rate and its ideal expenditure rate needed to meet consumption demands. The rate of change of labor allocation ($d a_j / dt$) is proportional to this income error. In vector notation, it's expressed as: $\dot{\mathbf{a}} = \Psi[\gamma M\mathbf{b} - N(\mathbf{p} \cdot \mathbf{c})\mathbf{a}]$ where $\mathbf{a}$ is the labor allocation vector, $\Psi$ is a reaction coefficient, $\gamma M$ is the mean money velocity, $\mathbf{b}$ is the money allocation vector, $N$ is the number of actors, $\mathbf{p}$ is the average price vector, and $\mathbf{c}$ is the consumption vector.

\item Money Equation: This equation describes how the allocation of money to different sectors changes based on the over- or under-production of commodities. It's based on the idea that commodities in oversupply will have lower average prices, leading to a decrease in that sector's income, and vice-versa for under-supplied commodities. The rate of change of sector income ($d b_j / dt$) is negatively proportional to the "production error" ($\epsilon_j(t)$), which is the difference between supply and demand for a commodity. In vector notation, it's expressed as: $\dot{\mathbf{b}} = -\omega N(\mathbf{A}\mathbf{l} - \mathbf{c})$ where $\omega$ is a reaction coefficient, $\mathbf{A}$ is a diagonal matrix of labor allocation proportions, $\mathbf{l}$ is the labor value vector, and $\mathbf{c}$ is the consumption vector.
\item Equilibrium Point and Stability: The paper shows that by setting the rates of change in both the labor and money equations to zero ($\dot{\mathbf{a}} = \mathbf{0}$ and $\dot{\mathbf{b}} = \mathbf{0}$), a unique equilibrium point can be found. At this equilibrium, the proportion of actors in a sector equals the proportion of money received by that sector, and this proportion is equal to the ratio of labor value to consumption rate ($l_j/c_j$). The paper further demonstrates the global asymptotic stability of this equilibrium point using Lyapunov's direct method, implying that the system will always approach this equilibrium regardless of its initial conditions.
\item  The Law of Value as an Attractor: A crucial result is Theorem 1, which states that labor values are global attractors for average market prices. This means that in equilibrium, the average price of a commodity ($\langle p_j \rangle$) is proportional to the labor-time required to make it ($l_j$), with a constant of proportionality ($\lambda$) known as the Monetary Expression of Labour-time (MELT).  $\lim_{t \to \infty} \mathbf{p}(t) = \lambda \mathbf{v}$ where $\mathbf{v}$ is the labor values vector.
    
\end{itemize}

Labour Commanded: The concept of "labour commanded" ($\kappa_j$) is introduced as the money price of a commodity divided by the MELT. This measures how much social labor-time a commodity fetches in the market. Deviations between labor commanded and embodied labor ($l_j$) serve as "labour reallocation signals," indicating over or under-valuation and guiding the adjustment of the division of labor.

In essence, the mathematical model formalizes how the dynamic interactions between labor allocation, money flows, and market prices lead to an equilibrium where prices reflect labor values and social labor is efficiently distributed, even in an economy with subjective pricing decisions.

\subsubsection{Labor Equation}

In a certain sense, this work continues the discussion started in the previous section, but now instead of an agent-based model, we have an equation-based model. We will therefore continue using the mathematical notations introduced earlier and define a few new ones.

\begin{itemize}
    \item The velocity of money is given by $\gamma M$, where $\gamma$ is a constant between 0 and 1.
\item The vector $\boldsymbol{b}\left(t\right)=\left(b_{1},\dots,b_{L}\right)$ is such that each element $b_j$ represents the instantaneous fraction (that is, at time $t$) of the total money flow being received by sector $j$.
\item Evidently, if we sum over all sectors we must have $\sum b_{j}=1$.
\end{itemize}


If the velocity of money $\gamma M$ is the total amount of money in circulation at a given instant $t$, and $b_j$ is the fraction of that amount that sector $j$ is receiving, then clearly the total amount of money that sector $j$ is receiving is given by the product $b_j \gamma M$.


If the velocity of money $\gamma M$ is the total amount of money in circulation at a given instant $t$, and $b_j$ is the fraction of that amount in sector $j$, then clearly the total amount of money that sector $j$ is receiving is given by the product $b_j \gamma M$.

The vector $\boldsymbol{a}\left(t\right)=\left(a_{1},\dots,a_{L}\right)$ is defined such that each element $a_j$ represents the current fraction of workers allocated in sector $j$.

We will use the average price $\left\langle p_{j}\right\rangle$ of the commodity to approximate the price distribution. Approximating that at each step, each agent consumes $1/c_j$ of commodity $j$, and that its price is $\left\langle p_{j}\right\rangle$, then the current price cost of consumption per step per agent is $\sum\left\langle p_{j}\right\rangle /c_{j}$.

Let us consider that agents switch sectors based on a signal indicated by prices. Each sector has an ideal expenditure rate that represents the amount of money that needs to be spent for the agents employed in this sector to satisfy their needs. This is given by the product of the individual consumption calculated previously and the current number of agents working in the sector ($a_j N$), that is, $a_j N \sum\left\langle p_{j}\right\rangle /c_{j}$.

Thus, the flow, or the error in the sector’s yield rate, is given by the difference between the total money flowing into sector $j$ and the ideal consumption:


$$\phi_{j}\left(t\right)=b_{j}\gamma M-a_{j}N\sum_{k=1}^{L}\frac{\left\langle p_{k}\right\rangle} {c_{k}}$$
Evidently, $\phi_{j}>0$ implies that the sector has profit, that is, the sector receives more money than the agents working in it need to satisfy their consumption. On the other hand, $\phi_{j} < 0$ implies a deficit, and $\phi_{j}=0$ represents the ideal equilibrium. We approximate that the switching of agents between sectors ($da_j/dt$) is proportional to $\phi_{j}$. For example, if there is a deficit, it means the sector does not have sufficient yield, so agents will leave the sector in search of another sector that is profitable and capable of absorbing them. In other words, $da_j/dt=\psi \phi_j(t)$, where $\psi$ is a positive constant that regulates the intensity of sector switching.

This is the \textbf{labor equation}: it defines how the allocation of labor across different production sectors varies according to the yield received by the sector from the sale of commodities. We can write it in a more complete form:


$$\frac{da_{j}}{dt}=\psi\phi_{j}\left(t\right)=\psi\left(b_{j}\gamma M-a_{j}N\sum_{k=1}^{L}\frac{\left\langle p_{k}\right\rangle}{c_{k}}\right)$$

Recalling that the consumption vector is given by $\boldsymbol{c}=(1/c_{1},\ldots,1/c_{L})$, and that we previously defined a price vector $\boldsymbol{p}=(p_{1},\ldots,p_{L})$, then $\boldsymbol{c}\cdot\boldsymbol{p}=\sum\frac{\left\langle p_{k}\right\rangle }{c_{k}}$.

Thus, the labor equation for the entire economy can be written in vector notation as:


$$\frac{d\boldsymbol{a}}{dt}=\psi\left(\gamma M\boldsymbol{b}-N\left(\boldsymbol{p}\cdot\boldsymbol{c}\right)\boldsymbol{a}\right)$$

Now, since an agent in sector $j$ produces $1/l_j$ units of the commodity at each step (it is produced after $l_j$ steps), the total production in the sector is given by $a_j N/l_j$.

If we define the average price of commodity $j$ as the result of dividing the sector’s income (as we saw earlier, given by $b_j \gamma M$) by the total quantity of commodities produced, that is, the money that entered the sector from the sale of commodities divided by the quantity of commodities produced, then the average price is:


$$\left\langle p_{j}\right\rangle =\frac{b_{j}\gamma M}{\frac{a_{j}N}{l_{j}}}=\frac{b_{j}\gamma Ml_{j}}{a_{j}N}=\frac{\gamma M}{N}\frac{b_{j}}{a_{j}}l_{j}$$

\subsubsection{Money Equation}

We now want to understand how the change in income depends on the distribution of labor. Certainly, a sector’s income depends on the quantity of commodities produced. Since no one consumes more than they need, the maximum consumption corresponds to the “social demand” for each commodity $j$.

As each agent consumes $1/c_j$ per step (meaning they consume one unit of the commodity every $c_j$ steps), the social demand is given by $N/c_j$. Similarly, each agent employed in sector $j$ produces $1/l_j$, and since the number of agents in the sector is $a_j N$, the total amount of commodity produced in the sector is $a_j N/l_j$.

Thus, analogous to what we did previously, the flow, or the production error given by the difference between the quantity of commodity produced and consumed, is given by:


$$\xi_{j}\left(t\right)=\frac{a_{j}N}{l_{j}}-\frac{N}{c_{j}}$$

Evidently, $\xi_{j}>0$ represents overproduction, and $\xi_{j}< 0$ implies that not enough is being produced, while $\xi_{j}=0$ denotes equilibrium. If we assume that when there is overproduction (underproduction) the price falls (rises), then the change in sector income is proportional, but inversely, to $\xi_{j}$, that is, $db_j/dt=-\omega \xi_j(t)$, where $\omega$ is a positive constant that adjusts the intensity with which the income distribution reacts to the production error.

This equation is known as the **money equation** because it defines how the allocation of money across different sectors varies according to overproduction or underproduction of commodities. It can be written in a more complete form as:


$$\frac{db_{j}}{dt}=-\omega\xi_{j}\left(t\right)=-\omega N\left(\frac{a_{j}}{l_{j}}-\frac{1}{c_{j}}\right)$$

If we define a matrix $A$ where the only nonzero elements are the diagonal elements $a_{ii}$, and using the production vector $\boldsymbol{l}=(1/l_1,...,1/l_L)$ as a column vector, we can write the money equation for the entire economy as:


$$\frac{d\boldsymbol{b}}{dt}=-N\omega\left(\boldsymbol{A}\boldsymbol{l}-\boldsymbol{c}\right)$$

Since the consumption vector is $\boldsymbol{c}=(1/c_1,...,1/c_L)$, we will now investigate how the equilibrium of this system, formed by the labor and money equations, is achieved.

\subsubsection{Equilibrium}

Bringing together the previous results, we have that a simple commodity system is described by the following system of $2L$ coupled differential equations (the labor equation and the money equation for the entire economy):


\begin{align*}
\dot{\boldsymbol{a}} & =\psi\left(\gamma M\boldsymbol{b}-N\left(\boldsymbol{p}\cdot\boldsymbol{c}\right)\boldsymbol{a}\right)\\
\dot{\boldsymbol{b}} & =-\omega N\left(\boldsymbol{A}\boldsymbol{l}-\boldsymbol{c}\right)
\end{align*}

Where:

$$\left\langle p_{j}\right\rangle =\frac{\gamma M}{N}\frac{b_{j}}{a_{j}}l_{j}$$


And with the following constraints:


\begin{align*}
\sum_{j=1}^{L}a_{j} & =1 & 0\leq a_{j}\leq1\\
\sum_{j=1}^{L}b_{j} & =1 & 0\leq b_{j}\leq1\\
\sum_{j=1}^{L}\frac{l_{j}}{c_{j}} & =1=\eta & l_{j},c_{j}>0\\
M,N & >0 & \omega,\psi>0\\
0\leq\gamma & \leq1
\end{align*}

And:  
- $M$ ($N$) is the total amount of money (population) in the system.  
- $a_{j}$ is the fraction of workers producing commodity $j$ at a given moment in time. $a_{j}N$ is then the number of workers in sector $j$, and $\boldsymbol{a}$ is the vector formed by the elements $a_{j}$. Meanwhile, $\boldsymbol{A}$ is a matrix where the only nonzero elements are the diagonal elements $(i,i)$, which are the elements $a_{i}$.  
- $\gamma$ is the fraction of money exchanged per unit of time, so $\gamma M$ is the velocity of money.  
- $b_{j}$ is the fraction of $\gamma$ received by the sector producing commodity $j$ at a given time, and $\boldsymbol{b}$ is the vector formed by the elements $b_{j}$.  
- $l_{j}$ ($c_{j}$) is the time required for a commodity $j$ to be produced (consumed). $\boldsymbol{l}$ ($\boldsymbol{c}$) is the vector formed by the elements $1/l_{j}$ ($1/c_{j}$).  
- $\omega$ ($\psi$) is a constant that regulates the intensity with which workers (income) migrate between sectors.  
- $\left\langle p_{j}\right\rangle$ is the average price of commodity $j$, and $\boldsymbol{p}$ is the vector formed by the elements $\left\langle p_{j}\right\rangle$.

In summary, this system works according to the following scheme: a certain division of labor results in over or underproduction of commodities, which activates a price correction mechanism based on supply and demand. This causes a change in sector income, which consequently leads agents to switch sectors, resulting in a new division of labor.  

This system has a single equilibrium point given by:


$$\boldsymbol{a}^{*}=\left(\frac{l_{1}}{c_{1}},\dots,\frac{l_{L}}{c_{L}}\right)=\boldsymbol{b}^{*}$$

To facilitate understanding, let us first revisit the equations for a single sector $j$:

\begin{align*}
\dot{a}_{j} & =\psi\left(b_{j}\gamma M-a_{j}N\sum_{k=1}^{L}\frac{\left\langle p_{k}\right\rangle }{c_{k}}\right)\\
\dot{b}_{j} & =-\omega N\left(\frac{a_{j}}{l_{j}}-\frac{1}{c_{j}}\right)
\end{align*}

Considering the equilibrium situation, that is, when there are no more variations ($\dot{a}_{j}=\dot{b}_{j}=0$):

\begin{align*}
0 & =\psi\left(b_{j}^{*}\gamma M-a_{j}^{*}N\sum_{k=1}^{L}\frac{\left\langle p_{k}^{*}\right\rangle }{c_{k}}\right)\\
0 & =-\omega N\left(\frac{a_{j}^{*}}{l_{j}}-\frac{1}{c_{j}}\right)
\end{align*}

Or simply:

\begin{align*}
0 & =b_{j}^{*}\gamma M-a_{j}^{*}N\sum_{k=1}^{L}\frac{\left\langle p_{k}^{*}\right\rangle }{c_{k}}\\
0 & =\frac{a_{j}^{*}}{l_{j}}-\frac{1}{c_{j}}
\end{align*}

From the second equation, we directly obtain:


$$a_{j}^{*}=\frac{l_{j}}{c_{j}}$$

Substituting into the first equation, and also replacing the definition of the average price, we have:

\begin{align*}
0 & =b_{j}^{*}\gamma M-\left(a_{j}^{*}\right)N\sum_{k=1}^{L}\frac{1}{c_{k}}\left(\left\langle p_{k}^{*}\right\rangle \right)\\
-b_{j}^{*}\gamma M & =-\left(\frac{l_{j}}{c_{j}}\right)N\sum_{k=1}^{L}\frac{1}{c_{k}}\left(\frac{\gamma M}{N}\frac{b_{k}^{*}}{a_{k}^{*}}l_{k}\right)\\
-b_{j}^{*} & =-\frac{l_{j}}{c_{j}}\sum_{k=1}^{L}\frac{l_{k}}{c_{k}}\frac{b_{k}^{*}}{a_{k}^{*}}
\end{align*}

Substituting again $a_{j}^{*}=\frac{l_{j}}{c_{j}}$ (a term that reappeared because of $\left\langle p_{k}\right\rangle$):


\begin{align*}
-b_{j}^{*} & =-\frac{l_{j}}{c_{j}}\sum_{k=1}^{L}\frac{l_{k}}{c_{k}}\frac{1}{a_{k}^{*}}b_{k}^{*}\\
-b_{j}^{*} & =-\frac{l_{j}}{c_{j}}\sum_{k=1}^{L}\frac{l_{k}}{c_{k}}\frac{c_{k}}{l_{k}}b_{k}^{*}\\
-b_{j}^{*} & =-\frac{l_{j}}{c_{j}}\sum_{k=1}^{L}b_{k}^{*}
\end{align*}

And since $\sum_{k=1}^{L}b_{k}^{*}=1$, then:

$$b_{j}^{*}=\frac{l_{j}}{c_{j}}=a_{j}^{*}$$

That is, when the system is in equilibrium, the proportion of agents employed in each sector is equal to the proportion of money each sector receives, which is proportional to $l_{j}/c_{j}$. This makes sense because all agents in this model require the same income to satisfy their needs. We had already shown that the efficient division of labor occurs when $a_{j}=l_{j}/c_{j}$, so a result we obtain without much effort is that, at equilibrium, the division of labor tends to maximum efficiency.

Furthermore, this equilibrium point is asymptotically stable, that is, regardless of the initial conditions, the system always evolves toward the equilibrium point. To prove this, we first need to a) linearize the system and b) shift the equilibrium point to the origin, then we can apply \href{https://eng.libretexts.org/Bookshelves/Industrial_and_Systems_Engineering/Book\%3A_Dynamic_Systems_and_Control_\%28Dahleh_Dahleh_and_Verghese\%29/13\%3A_Internal_\%28Lyapunov\%29_Stability/13.03\%3A_Lyapunov\%27s_Direct_Method}{Lyapunov's direct method}. 

The equations are:


\begin{align*}
\dot{a}_{j} & =\psi\gamma Mb_{j}-\psi\sum_{k=1}^{L}\frac{l_{k}\gamma M}{c_{k}}\frac{a_{j}b_{k}}{a_{k}}\\
\dot{b}_{j} & =\frac{\omega N}{c_{j}}-\frac{\omega N}{l_{j}}a_{j}
\end{align*}

We need to linearize the first term. First, since $\sum_{j=1}^{L}a_{j}=1$, it follows that $\sum_{j=1}^{L}\dot{a}_{j}=0$, because the variation in each element must be 'cancelled' by an equivalent variation of opposite sign in another (or other) elements so that the sum remains constant. Thus:


\begin{align*}
\dot{a}_{j} & =\psi\gamma Mb_{j}-\psi\sum_{k=1}^{L}\frac{l_{k}\gamma M}{c_{k}}\frac{a_{j}b_{k}}{a_{k}}\\
\sum_{j=1}^{L}\dot{a}_{j} & =\sum_{j=1}^{L}\left[\psi\gamma Mb_{j}-\psi\sum_{k=1}^{L}\frac{l_{k}\gamma M}{c_{k}}\frac{a_{j}b_{k}}{a_{k}}\right]\\
0 & =\sum_{j=1}^{L} b_{j}-\sum_{j,k=1}^{L}\frac{l_{k}}{c_{k}}\frac{a_{j}b_{k}}{a_{k}}\\
\sum_{j=1}^{L}b_{j} & =\sum_{j,k=1}^{L}\frac{l_{j}}{c_{k}}\frac{a_{j}b_{k}}{a_{k}}\\
\left(\sum_{j=1}^{L}b_{j}\right) & =\left(\sum_{j=1}^{L}a_{j}\right)\left(\sum_{k=1}^{L}\frac{l_{k}}{c_{k}}\frac{b_{k}}{a_{k}}\right)\\
1 & =\sum_{k=1}^{L}\frac{l_{k}}{c_{k}}\frac{b_{k}}{a_{k}}
\end{align*}

Where we use the fact that $\sum_{j=1}^{L}b_{j}=\sum_{j=1}^{L}a_{j}=1$.  
Revisiting our equation then:


\begin{align*}
\dot{a}_{j} & =\psi\gamma Mb_{j}-\psi\sum_{k=1}^{L}\frac{l_{k}\gamma M}{c_{k}}\frac{a_{j}b_{k}}{a_{k}}\\
 & =\psi\gamma Mb_{j}-\psi\gamma Ma_{j}\sum_{k=1}^{L}\frac{l_{k}}{c_{k}}\frac{b_{k}}{a_{k}}\\
 & =\psi\gamma Mb_{j}-\psi\gamma Ma_{j}\psi
\end{align*}

That is:

$$\dot{a}_{j}=\psi\gamma M\left(b_{j}-a_{j}\right)$$


Now, with a simple change of variables $x_{j}=a_{j}-\frac{l_{j}}{c_{j}}$ and $y_{j}=b_{j}-\frac{l_{j}}{c_{j}}$, we shift the equilibrium point to the origin. With this change, for the labor equation, we have:


\begin{align*}
\dot{a}_{j} & =\psi\gamma M\left(b_{j}-a_{j}\right)\\
\frac{d}{dt}\left(x_{j}+\frac{l_{j}}{c_{j}}\right) & =\psi\gamma M\left(\left[y_{j}+\frac{l_{j}}{c_{j}}\right]-\left[x_{j}+\frac{l_{j}}{c_{j}}\right]\right)\\
\dot{x}_{j} & =\psi\gamma M\left(y_{j}-x_{j}\right)
\end{align*}

Ou em notação vetorial:

$$
\dot{\boldsymbol{x}}=\psi\gamma M\left(\boldsymbol{y}-\boldsymbol{x}\right)
$$

And for the money equation:

\begin{align*}
\dot{b}_{j} & =-\omega N\left(\frac{a_{j}}{l_{j}}-\frac{1}{c_{j}}\right)\\
\frac{d}{dt}\left(y_{j}+\frac{l_{j}}{c_{j}}\right) & =-\omega N\left(\frac{1}{l_{j}}\left(x_{j}+\frac{l_{j}}{c_{j}}\right)-\frac{1}{c_{j}}\right)\\
\dot{y}_{j} & =-\omega N\left(\frac{x_{j}}{l_{j}}\right)
\end{align*}

Or in vector notation:

$$
\boldsymbol{y}=-\omega n\boldsymbol{X}\boldsymbol{l}
$$

We then have:


\begin{align*}
\dot{x}_{j} & =\psi\gamma M\left(y_{j}-x_{j}\right)\\
\dot{y}_{j} & =-\omega N\left(\frac{x_{j}}{l_{j}}\right)
\end{align*}

Or:

\begin{align*}
\dot{\boldsymbol{x}} & =\psi\gamma M\left(\boldsymbol{y}-\boldsymbol{x}\right)\\
\boldsymbol{y} & =-\omega n\boldsymbol{X}\boldsymbol{l}
\end{align*}

Analogously to $\boldsymbol{A}$, $\boldsymbol{X}$ is the matrix where only the diagonal elements $x_{ii}$ are nonzero. Since the equilibrium is now at the origin, the terms $x_{j}$ and $y_{j}$ represent an error in production and income in sector $j$.  

To formalize this, we define a function $V:\mathbb{R}^{2L}\rightarrow\mathbb{R}$, meaning that $V$ assigns a real number to each vector of variables $\left(x_{1},\dots,x_{L},y_{1},\dots,y_{L}\right)$. This function allows us to assign a scalar value to each configuration of the system. More explicitly:

$$
V\left(x_{1},\dots,x_{L},y_{1},\dots,y_{L}\right)=\frac{1}{2\psi\gamma M}\sum x_{j}^{2}+\frac{1}{2\omega N}\sum l_{j}y_{j}^{2}
$$

We thus have a scalar measure of the error for each possible state of our system. Obviously, if we are at the origin $x_{j}=y_{j}=0$, then $V=0$, and if we are away from the origin, then $V>0$, since $x_{j}$ and $y_{j}$ are real numbers greater than or equal to zero. All terms are necessarily non-negative, and the choice of this specific form for $V$ is to simplify later calculations.  

What we want is for $V$ to be a Lyapunov function. For this to be possible, $V$ must first satisfy $V=0$ at the origin and $V>0$ in the neighborhood of the origin. As discussed, both cases are satisfied with this definition of $V$. The last requirement is that $\dot{V}\leq0$. Taking the derivative:


\begin{align*}
\frac{d}{dt}V= & \frac{d}{dt}\left(\frac{1}{2\psi\gamma M}\sum x_{j}^{2}+\frac{1}{2\omega N}\sum l_{j}y_{j}^{2}\right)\\
= & \frac{1}{2\psi\gamma M}\sum\frac{d}{dt}\left(x_{j}^{2}\right)+\frac{1}{2\omega N}\sum l_{j}\frac{d}{dt}\left(y_{j}^{2}\right)\\
= & \frac{1}{2\psi\gamma M}\sum2x_{j}\dot{x}_{j}+\frac{1}{2\omega N}\sum l_{j}2y_{j}\dot{y}_{j}\\
\dot{V}= & \frac{1}{\psi\gamma M}\sum x_{j}\dot{x}_{j}+\frac{1}{\omega N}\sum l_{j}y_{j}\dot{y}_{j}
\end{align*}

Since the derivative of a sum is the sum of the derivatives, substituting $\dot{x}_{j}$ and $\dot{y}_{j}$:

\begin{align*}
\dot{V}= & \frac{1}{\psi\gamma M}\sum x_{j}\dot{x}_{j}+\frac{1}{\omega N}\sum l_{j}y_{j}\dot{y}_{j}\\
= & \frac{1}{\psi\gamma M}\sum x_{j}\left[\psi\gamma M\left(y_{j}-x_{j}\right)\right]+\frac{1}{\omega N}\sum l_{j}y_{j}\left[-\omega N\left(\frac{x_{j}}{l_{j}}\right)\right]\\
= & \sum x_{j}\left(y_{j}-x_{j}\right)-\sum y_{j}x_{j}\\
= & \sum x_{j}y_{j}-\sum x_{j}^{2}-\sum y_{j}x_{j}
\end{align*}

Or simply:


$$
\dot{V}=-\sum x_{j}^{2}
$$

Evidently, $\dot{V}\leq0$, as we wanted. Moreover, $\dot{V}=0$ occurs only for $\boldsymbol{x}^{*}=0$; in all other cases, we have $\dot{V}< 0$. That is, the scalar error given by $V$ always decreases over time until equilibrium is reached, where the error is zero and then stabilizes. $\dot{V}< 0$ in the neighborhood of the equilibrium point is what defines a point as asymptotically stable. In this case, since the Lyapunov function is valid for all of $\mathbb{R}^{2L}$ and not just near the equilibrium point, we can say that the point is globally asymptotically stable.

Finally, we can reach the result that interests us most: the **law of value**.  Labor values are attractors for the average prices. That is:


$$\lim_{t\rightarrow\infty}\boldsymbol{p}\left(t\right)=\lambda\boldsymbol{v}$$

Where $\boldsymbol{v}=\left(l_{1},\dots,l_{L}\right)$. Since $1/l_{j}$ denotes how many units of commodity $j$ are produced per unit of time, $l_{j}$ denotes how many units of time are required to produce commodity $j$. Returning to the definition of the average price and using the equilibrium point ($a_{j}=b_{j}$), we have:


$$
\left\langle p_{j}\right\rangle =\frac{\gamma M}{N}\frac{b_{j}}{a_{j}}l_{j}=\left(\frac{\gamma M}{N}\right)l_{j}
$$

Where evidently we then have $\lambda=\gamma M/N$, representing the monetary value of one unit of labor-time. As the system tends toward equilibrium, the prices $\boldsymbol{p}$ tend toward $\lambda\boldsymbol{v}$.

It is interesting to note that the numerator of $\lambda$ is the same as that calculated for the simulation: the velocity of money. But the denominator is different. We must remember that the denominator is the labor expended and exchanged in the market over a given time interval. If our time interval is one unit, and we have $N$ agents in a system where all agents are always necessarily producing, then in one unit of time, $N$ units of labor-time are incorporated into the commodities.

Since we are analyzing only the goods being exchanged in the market, we can see on average, $N$ units of labor-time are exchanged per unit of time. We can also note that, in a way, when agents switch sectors seeking to maximize their income, they always try to work in a sector where the commodity they produce will be better sold.


\subsubsection{Conclusion}


The quantitative part of Marxist value theory tells us that prices are related to the socially necessary labor time required to produce a commodity, and deviations of prices from values are signals of social errors that serve to redistribute labor. To analyze this more closely, let us define the concept of labour commanded: $\left\langle \epsilon_{j}\right\rangle=\frac{\left\langle p_{j}\right\rangle }{\lambda}$.

This quantity represents how much socially necessary labor time, on average, a commodity “contains” when it is exchanged in the market. Recall that the MELT, or “Monetary Expression of Labor Time,” has units of money per unit of labor-time. Since the price is in money, $\epsilon$ has units of labor-time. If $\left\langle \epsilon_{j}\right\rangle < l_j$, this means the commodity has been traded at a depreciated price; obviously, if $\left\langle \epsilon_{j}\right\rangle > l_j$, then the commodity is overvalued. Evidently, all these quantities are objective and independent of any individual subjectivity about the utility of the commodity, and at equilibrium $\left\langle \epsilon_{j}\right\rangle = l_j$.  

Manipulating the necessary labor, it can be rewritten as:


$$\left\langle \epsilon_{j}\right\rangle =\frac{1}{\lambda}\left\langle p_{j}\right\rangle =\frac{N}{\gamma M}\frac{\gamma M}{N}\frac{b_{j}}{a_{j}}l_{j}=\frac{b_{j}}{a_{j}}l_{j}
$$

That is, $b_j=\left\langle \epsilon_{j}\right\rangle \frac{a_{j}}{l_{j}}$. Substituting into the linearized version of the labor equation, we have:


$$\dot{a}_{j}=\psi\gamma M\left(b_{j}-a_{j}\right)=\psi\gamma Ma_{j}\left(\frac{\left\langle \epsilon_{j}\right\rangle }{l_{j}}-1\right)$$

Or simply:


$$\dot{\boldsymbol{a}}=\psi\gamma M\boldsymbol{a}\left(\frac{\left\langle \epsilon_{j}\right\rangle }{l_{j}}-1\right)$$

Looking at the term in parentheses, if $\left(\frac{\left\langle \epsilon_{j}\right\rangle }{l_{j}}>1\right)$, then the subtraction is positive and the commodity is overvalued, signaling the need to expand the sector and increase its production. Evidently, $\left(\frac{\left\langle \epsilon_{j}\right\rangle }{l_{j}}< 1\right)$ signals the opposite. It is worth noting that the equilibrium we are working with is a statistical equilibrium, meaning we are discussing an ‘average’ relationship between value and price. Even at equilibrium, an individual transaction can occur at different prices. The law of value, as presented, states that at equilibrium, the average price of commodities is proportional to their labor values.

It is important to note that this model does not exhaust the possibilities for investigating the topic. There are models that explore, for example, complex production systems where one commodity is required to produce another (Krause), but I believe that for a first contact, the content presented here already fulfills its purpose.  

If we define a coefficient $\alpha_{ij}$, which indicates that 1 hour of labor in the production of commodity $i$ is equivalent to $\alpha_{ij}$ hours of labor in commodity $j$, we can then write:


$$
\alpha_{ij}=\frac{\left\langle p_{i}\right\rangle /l_{i}}{\left\langle p_{j}\right\rangle /l_{j}}=\frac{\frac{\gamma M}{N}\frac{b_{i}}{a_{i}}l_{i}/l_{i}}{\frac{\gamma M}{N}\frac{b_{j}}{a_{j}}l_{j}/l_{j}}=\frac{b_{i}/a_{i}}{b_{j}/a_{j}}
$$

First of all, we can note that $\frac{\left\langle p_{i}\right\rangle }{l_{i}}$ gives us the amount of money that one unit of labor-time in the production of commodity $i$ is measured in money. So, when we divide this amount for commodity $i$ by the corresponding amount for commodity $j$, we obtain the “price” of one unit of labor-time in the production of commodity $i$ expressed in units of labor-time for the production of commodity $j$.  

In other words, this ratio “indicates that 1 hour of labor in the production of commodity $i$ is equivalent to $\alpha_{ij}$ hours of labor in commodity $j$.” For example, if commodity $i$ has 2 units of money per unit of labor-time and commodity $j$ has 4, then $\alpha_{ij}=1/2$, meaning that for each time interval in which commodity $j$ is equivalent to 1 unit of money, this time corresponds to 0.5 units of money for commodity $i$.

We have already obtained the result that at equilibrium $a_j=b_j$, so evidently at equilibrium $b_j/a_j=1$, and consequently $\alpha_{ij}=1$. In other words, the statement that values are attractors for prices at equilibrium is equivalent to the statement that homogeneous coefficients are attractors for heterogeneous coefficients at equilibrium.  

The law of value is a dynamic theory of labor allocation based on the tendency of heterogeneous labor to be homogenized through commodity exchange. Evidently, this is an idealized scenario that allows us to understand some fundamental concepts for analyzing capitalism and is far from exhausting the subject.

In conclusion, the law of value is a phenomenon that emerges from the dynamic interaction of private commodity producers. In the model presented: values are global attractors for prices, prices in turn act as signals of errors that work to allocate the social division of available labor more efficiently across sectors, and the tendency of prices to approach value is the monetary expression of the tendency to allocate social labor efficiently. The MELT is an interesting constant because it summarizes a non-obvious relationship between a production phenomenon (production time) and a market phenomenon (prices). With the definition of labour commanded, we can have an indirect signal of whether the labor in the production of that commodity is socially necessary or not.

Perhaps most interestingly, the law of value operates “behind the scenes”; it emerges naturally from the system dynamics and does not need to be imposed directly at any point. It is also worth noting that the system’s equilibrium is statistical, meaning we look at the average price of a commodity, which can individually be exchanged at different prices even when the system reaches equilibrium. Moreover, the law of value can only emerge in an economic system that contains a feedback cycle between production, consumption, exchange, and the reallocation of labor resources.
